\documentclass[12pt,nonacm,acmlarge,titlepage,screen]{acmart}

% make ACM settings and load extra packages
% settings specific to the ACM class

%% DO NOT MODIFY THIS FILE UNLESS IT IS FOR SUBMITTING TO AN ACM CONFERENCE/JOURNAL

%% Rights management information.  This information is sent to you
%% when you complete the rights form.  These commands have SAMPLE
%% values in them; it is your responsibility as an author to replace
%% the commands and values with those provided to you when you
%% complete the rights form.
% \setcopyright{acmlicensed}
\setcopyright{none}
% \copyrightyear{2018}
% \acmYear{2018}
\acmDOI{XXXXXXX.XXXXXXX}

%% These commands are for a JOURNAL article.
% \acmJournal{POMACS}
% \acmVolume{37}
% \acmNumber{4}
% \acmArticle{111}
% \acmMonth{8}

%% Submission ID.
%% Use this when submitting an article to a sponsored event. You'll
%% receive a unique submission ID from the organizers
%% of the event, and this ID should be used as the parameter to this command.
% \acmSubmissionID{123-A56-BU3}


%% citation style
% \citestyle{acmauthoryear}
\citestyle{acmnumeric}


%% ACM dates
% \received{\docDate}
% \received[revised]{12 March 2009}
% \received[accepted]{5 June 2009}
%% This file contains extra packages to use on top of the ACM template

%% STYLING ToC
\usepackage{titletoc}

%% SUB-FIGURES
\usepackage{subcaption}

%% ALGORITHMS
\usepackage{algorithm}
\usepackage{algorithmicx}
\usepackage{algpseudocode}

%% MORE EQUATIONS TOOLS
\usepackage{amsmath,mathtools}

%% CODE LISTINGS
\usepackage{listings}

% EASIER CROSS-REFERENCES
\usepackage[capitalise,nameinlink,noabbrev]{cleveref}

% TO GET ORDINAL NUMBERS
\usepackage[super]{nth}

% PAGE-BREAKABLE TABLES
\usepackage{longtable}

% IN-TABLE TEXT ALIGNMENTS
\usepackage{array}

% document details
% replace those with your details
\def\docTitle{A Grounded University Support Assistant: Retrieval-Augmented Generation with Confidence-Based Escalation}
\def\docAuthor{\textsc{Muhammad Faizal Haris Aidi}}
\def\docAuthorID{H00459013}
\def\docSupervisor{\textsc{Dr. Rosalind Deena Kumari}}
\def\docDegree{BSc (Hons) Computing Science}
%\def\docDegree{MSc Cookery}
\def\docType{Research Report}
% \def\docType{Honours Dissertation}
% \def\docType{Masters Dissertation}
\def\docCourse{F20PA/B/C - Honours Project}
% \def\docCourse{F21RP - Research Methods and Project Planning}
% \def\docCourse{F21MP - Masters Project and Dissertation}

% load extra formatting
% formatting and package settings


% document details
\def\docDate{\nth{\the\day}~\ifcase\the\month\or
  January\or February\or March\or April\or May\or June\or
  July\or August\or September\or October\or November\or
  December\fi~\the\year}
\def\docInstitution{Heriot-Watt University}
\def\docSchool{School of Mathematical and Computer Sciences}
\author{\docAuthor}
\title{\docTitle}
\date{\docDate}

% change geometry of ACM template
\geometry{twoside=true, head=13pt, foot=2pc,
     paperwidth=21cm, paperheight=29.7cm,
     includeheadfoot,
     top=78pt, bottom=114pt, inner=81pt, outer=81pt,
     marginparwidth=4pc,heightrounded
 }%

% to make sections start on new pages
% \AddToHook{cmd/section/before}{\newpage}

% colors
% main text
\def\textColor{\color[HTML]{212121}} 
\def\secTextColor{\color[HTML]{424242}}
% footer and header
\def\headerColor{\color[HTML]{565656}} 
% title, chapters, sections
\def\titleColor{\color[HTML]{212121}}

% set line spacing
% \onehalfspacing
% \singlespacing

% set table cells padding
\def\arraystretch{1.5}

% set headings format
% unset to follow ACM format
% \titleformat{\chapter}{\singlespacing\normalfont\huge\bfseries\titleColor}{\chaptertitlename\ \thechapter}{1em}{}
% \titleformat{\chapter}{\singlespacing\normalfont\huge\bfseries\titleColor}{\thechapter}{1em}{}
% \titleformat{\section}{\singlespacing\normalfont\large\bfseries\titleColor}{\thesection}{1em}{}
% \titleformat{\subsection}{\singlespacing\normalfont\large\bfseries}{\thesubsection}{1em}{}
% \titleformat{\subsubsection}{\singlespacing\normalfont\normalsize\bfseries}{\thesubsubsection}{1em}{}
% \titlespacing*{\chapter}{0pt}{0pt}{10pt}

% set paragraph indent
% unset to follow ACM format
% \setlength{\parindent}{18pt} 

% set the text colour
% \textColor 

% set hyperlinks + pdf metadata
\hypersetup{
    pdftitle={\docTitle},
    pdfauthor={\docAuthor},
    pdfkeywords={computer science dissertation},
    colorlinks=true,%
    citecolor=[HTML]{152845},%
    linkcolor=[HTML]{152845},%
    urlcolor=[HTML]{1976D2}
}

% styling the code listings
\definecolor{codegreen}{rgb}{0,0.6,0}
\definecolor{codegray}{rgb}{0.5,0.5,0.5}
\definecolor{codepurple}{rgb}{0.58,0,0.82}
\definecolor{backcolour}{rgb}{0.95,0.95,0.92}

\lstdefinestyle{mystyle}{
    backgroundcolor=\color{backcolour},   
    commentstyle=\color{codegreen},
    keywordstyle=\color{magenta},
    numberstyle=\tiny\color{codegray},
    stringstyle=\color{codepurple},
    basicstyle=\ttfamily\singlespacing,
    breakatwhitespace=false,         
    breaklines=true,                 
    captionpos=b,                    
    keepspaces=true,                 
    numbers=left,                    
    numbersep=5pt,                  
    showspaces=false,                
    showstringspaces=false,
    showtabs=false,                  
    tabsize=2
}
\lstset{style=mystyle}

% reset names for code listing
\renewcommand{\lstlistingname}{Code Listing}
\renewcommand{\lstlistlistingname}{List of \lstlistingname s}

% set cref for listing
\crefname{lstlisting}{\lstlistingname}{\lstlistingname s}
\Crefname{lstlisting}{\lstlistingname}{\lstlistingname s}

% set cross-references labels with cleveref
\crefname{appendixchapter}{Appendix}{Appendices}
\Crefname{appendixchapter}{Appendix}{Appendices}


% set custom headers/footers
% unset to follow ACM template
% \pagestyle{fancy}
% \fancyhf{}
% \renewcommand{\headrulewidth}{0pt}
% \renewcommand{\footrulewidth}{0pt}
% \fancyhead[R]{\headerColor \docType}
% \fancyhead[L]{\headerColor \docAuthor}
% \fancyfoot[C]{\headerColor \thepage}
% % Redefine the plain page style for chapter pages
% \fancypagestyle{plain}{%
%     \fancyhf{}%
%     \fancyfoot[C]{\headerColor \thepage}
% }

% custom section heading for front-matter content with addition to the ToC
\newcommand{\frontmattersection}[2]{
    \thispagestyle{plain} % remove the header on page
    % \vspace*{1cm} % to align vertically
    \phantomsection
    \begin{center} 
        \textsc{#1} 
        \addcontentsline{toc}{section}{#1}
    \end{center}
    #2 % Main Text
    \clearpage % leave blank space underneath
}

\renewcommand{\abstract}[1]{
    \frontmattersection{Abstract}{#1}
}

\newcommand{\acknowledgements}[1]{
    \frontmattersection{Acknowledgements}{#1}
}

% formatting the table of content
\dottedcontents{section}[1em]{\bfseries}{1.2em}{0pc}
\dottedcontents{subsection}[3em]{}{2.2em}{1pc}

\renewcommand{\contentsname}{Table of Contents}

\begin{document}

%%%%%%%%%%%%%%%%%%%%%%%%%%%%%%%%%%%%%%%%%%%%%%%%%%
% FRONT MATTER
%%%%%%%%%%%%%%%%%%%%%%%%%%%%%%%%%%%%%%%%%%%%%%%%%%
\begin{titlepage}

\begin{center}
	
	% Title (and subtitle)
	~\\[0.2cm]
	{\huge \titleColor \textbf{\docTitle}}\\
	~\\[0.7cm]
	
	% Author(s)
	{\Large \textit{by} \docAuthor}\\[0.7cm]
	{\large \docType}\\[1cm]
	{\large \textit{Submitted in partial fulfilment for the degree of}\\[0.2cm] \docDegree}\\[0.7cm]
	{\large \textit{Supervised by} {\Large \docSupervisor}}
	~\\[1cm]
	
	% Organisation	
	% ~\\[0cm]
        % logo / coat of arms
        \begin{figure}[h]
            \centering
		\includegraphics[height=3cm,keepaspectratio]{images/hwcoa_bw.jpg}
        \Description{Heriot-Watt University's coat-of-arms}
	\end{figure}
	\begin{large}
		\textsc{\docInstitution}\\
		\docSchool\\
% 		Department\\
	\end{large}
	~\\[1cm]
	
	% Date
	{\large \docDate}\\
	~\\[1.5cm]
	
	% Copyright
	\begin{minipage}{0.7\textwidth}
	% \singlespacing
	\begin{footnotesize}
	The copyright in this dissertation is owned by the author. Any quotation from the dissertation or use of any of the information contained in it must acknowledge it as the source of the quotation or information.
	\end{footnotesize}
	\end{minipage}
	
\end{center}

\end{titlepage}

\newpage
\pagenumbering{roman}

% authorship declaration
% this declaration is automatically filled with the details provided in main.tex

\thispagestyle{plain} % remove the header on page

\phantomsection
\begin{center} 
    \textsc{Declaration of Authorship} 
    \addcontentsline{toc}{section}{Declaration of Authorship}
\end{center}


\begin{center} 
    \begin{tabular}{|p{0.3\textwidth}|p{0.675\textwidth}|}
        \hline
        \textbf{Course code and name} & \docCourse \\
        \hline
        \textbf{Type of assessment} & Individual \\
        \hline
        \textbf{Document title} & \docTitle \\
        \hline
        \textbf{Student name} & \docAuthor \\
        \hline
        \textbf{Student ID} & \docAuthorID \\
        \hline
    \end{tabular}
\end{center}

\bigskip

\begin{longtable}{|>{\raggedright\arraybackslash}p{\textwidth}|}
    \hline
    \endfirsthead
    \hline
    \multicolumn{1}{|r|}{\footnotesize\it (form continued from previous page)} \\
    \endhead
    \multicolumn{1}{|r|}{\footnotesize\it (form continues on next page)} \\
    \hline
    \endfoot
    \hline
    \endlastfoot
    
        \textbf{Statement on Generative AI Use:}
        \\
        
        EDIT this statement to indicate what use you made of Generative AI for your project work and in preparation of this document. If your project is about developing or experimenting with Generative AI, you still need to indicate your Generative AI use beyond the developments and experiments. If you have not used Generative AI, indicate it explicitly.
        \\
    \hline
    \multicolumn{1}{l}{}
    \\
    \hline
        \textbf{By signing this form:}
        \\
        
        \textbullet~ I declare that the work I have submitted for this assessment is my own original work. I have not presented the ideas, words, data, code, designs, analyses, or other intellectual contributions of another person or system as my own. 
        \\
        
        \textbullet~ My submission is expressed in my own words except where the work of others has been appropriately acknowledged, cited, and referenced. All sources used in the preparation of this assessment, including books, journal articles, websites, software tools, and Generative AI tools, have been properly acknowledged.
        \\
        
        \textbullet~ Where Generative AI tools have been used, I have complied with the course’s stated policies regarding Generative AI use. I have clearly and accurately disclosed in the above statement box all uses of Generative AI in the assessment, including the specific tools used, the purpose for which they were used, and the extent to which any generated content influenced the final submission.
        \\
        
        \textbullet~ I understand that I remain fully responsible for the content, accuracy, integrity, and originality of my submission, regardless of any assistance received from Generative AI or other tools.
        \\

        \textbullet~ I confirm that I have read, understood and followed the University's Regulations on academic misconduct as published on the \href{https://www.hw.ac.uk/uk/services/academic-registry/academic-integrity/academic-misconduct.htm}{University's website}, and that I am aware of the penalties that I will face should I not adhere to the University Regulations.
        \\
        
        \textbullet~ I confirm that I have read, understood and avoided the different types of plagiarism explained in the University guidance on \href{https://heriotwatt.sharepoint.com/sites/skillshub/SitePages/Academic-Integrity-and-Plagiarism.aspx}{Academic Integrity and Plagiarism}.
        \\
        
        \\

        \noindent \textbf{Student signature}: \docAuthor
        \\
        
        \\
        
        \noindent \textbf{Date}: \docDate
        \\
    \hline
\end{longtable}

\noindent \textbf{Your work will not be marked if a signed copy of this form is not included with your submission.}



\clearpage

% Acknowledgements
\acknowledgements{
Acknowledgements text.
}

% Abstract
\abstract{
University students often struggle to find accurate answers about fees, deadlines and regulations that are spread across web pages, PDFs and emails, and general-purpose chatbots can answer such questions confidently yet incorrectly. This project designs, implements and evaluates a university support assistant that uses retrieval-augmented generation (RAG) to answer only from official university documents, cites its sources, and escalates to staff when its confidence is low. Four systems will be compared on the same knowledge base and a purpose-built test set of about 200 student questions: an FAQ-matching baseline, a plain large language model, RAG with citations, and RAG with confidence-based escalation. The evaluation will measure answer correctness, hallucination, citation accuracy, retrieval quality, escalation precision and recall, latency and cost, using paired statistical tests, and may include a small usability study subject to ethics approval. This report presents the motivation, related work, methodology and feasibility work that show the project can be completed in Semester 2.
}

% table of contents

\newpage

\thispagestyle{plain} % remove the header on page

\setcounter{tocdepth}{2} % set depth of ToC, 0 = cha, 1 = sec, etc.

\tableofcontents % the actual table of contents command

\clearpage

	
%%%%%%%%%%%%%%%%%%%%%%%%%%%%%%%%%%%%%%%%%%%%%%%%%%
% MAIN BODY
%%%%%%%%%%%%%%%%%%%%%%%%%%%%%%%%%%%%%%%%%%%%%%%%%%
\newpage
\pagenumbering{arabic}

% Introduction
%%%%%%%%%%%%%%%%%%%%
\section{Introduction} \label{cha:intro}

Students need timely and accurate answers about admissions, course registration, assessment deadlines, fees, financial aid and academic regulations. At Heriot-Watt University Malaysia this information is spread across web pages, PDF policies and emails, so students must search several places and staff answer many repetitive queries. Large language models (LLMs) can answer such questions fluently, but they can also produce confident statements that no source supports, a problem known as hallucination \cite{ji2023survey}. For university-specific facts such as a fee, a deadline or a late-submission penalty, an unsupported answer can cause real harm. This project designs, implements and evaluates a university support assistant that answers only from official documents, shows its sources, and hands the question to staff when it is not confident.

\subsection{Background and Motivation} \label{sec:intro_motiv}

Retrieval-augmented generation (RAG) \cite{lewis2020rag} addresses part of this problem. It retrieves relevant passages from a document collection and gives them to an LLM as context, so that answers can be grounded in, and traced back to, those passages. Surveys describe how RAG has grown into a family of designs for retrieval, generation and augmentation \cite{gao2023rag}. RAG alone does not guarantee a correct answer: retrieval can miss the relevant passage, and a generator may still answer from its own knowledge when the retrieved text is insufficient. A complementary line of work studies abstention, where a system declines to answer when it should not \cite{wen2025know}. Evaluating this behaviour needs questions that the knowledge base cannot answer, in addition to answerable ones \cite{peng2025unanswerability}.

University chatbots built on RAG already exist. Marcel, for example, answers admission questions from scraped university web pages, hides its sources when it abstains, and evaluates retrieval and generation on questions that are answerable and unanswerable from its knowledge base \cite{trienes2025marcel}. Earlier admissions chatbots used intent-based dialogue management instead \cite{nguyen2021neu}. In Marcel the decision to abstain is left to the answer-generation prompt, and the authors suggest that other ways of judging answerability are worth investigating \cite{trienes2025marcel}. This project builds on that observation. Rather than only building another chatbot, it treats the decision to answer or escalate as a measurable component, and compares it against simpler alternatives on a test set specific to the university.

\subsection{Aim and Objectives} \label{sec:intro_aim}

The aim is to design, implement and rigorously evaluate a grounded, source-citing, escalation-aware chatbot for university student support, and to quantify its benefits and limits against simpler alternatives. The objectives are:

\begin{description}
  \item[O1] Review the literature on educational chatbots, RAG, hallucination and abstention, and identify the gap this project addresses.
  \item[O2] Build a versioned knowledge base from public university documents.
  \item[O3] Implement four comparable systems: an FAQ-matching baseline, a plain LLM, RAG with citations, and RAG with confidence-based escalation.
  \item[O4] Build a labelled test set of about 200 student questions across topics and difficulty types, including questions that cannot be answered from the documents.
  \item[O5] Implement a confidence-based escalation mechanism and tune its threshold on a development split.
  \item[O6] Evaluate the systems on correctness, faithfulness, retrieval quality, escalation quality, latency and cost, using paired statistical tests.
  \item[O7] Deliver a web interface for demonstration and, if ethics approval is granted, a small usability study.
  \item[O8] Reflect critically on limitations, threats to validity and ethical issues.
\end{description}

The project investigates four research questions:

\begin{description}
  \item[RQ1] How much more accurate and faithful is RAG than an FAQ-matching baseline and a plain LLM on university-specific questions?
  \item[RQ2] Which retrieval design choices (chunking, number of retrieved passages, dense, lexical or hybrid retrieval) matter most?
  \item[RQ3] Does confidence-based escalation reduce incorrect answers, and what fraction of answerable questions does it wrongly escalate?
  \item[RQ4] (optional, subject to ethics approval) Do students find the assistant usable and trustworthy compared with searching the university website?
\end{description}

Meeting these objectives means solving four problems: turning heterogeneous web pages and PDFs into a clean, versioned knowledge base; generating answers that stay within the retrieved evidence and cite it; deciding when the evidence is too weak to answer; and measuring all of this fairly, including on questions that should not be answered.

\subsection{Organisation} \label{sec:intro_orga}

Section~\ref{cha:back} reviews related work on RAG, hallucination, abstention, university chatbots and the evaluation of RAG systems. The following sections set out the methodology and requirements, and the feasibility work carried out so far, including the pilot study, and then conclude. The back matter contains the references, the project management plan (Gantt chart, risk analysis and professional, legal, ethical and social issues), a summary of generative AI use, and the appendices.


% Background Research
%%%%%%%%%%%%%%%%%%%%
\section{Background Research} \label{cha:back}

This section reviews the work that the project builds on. It covers retrieval-augmented generation and the retrieval methods behind it, hallucination, abstention, existing university support chatbots, and how such systems are evaluated. It ends by naming the gaps that lead to the research questions in Section~\ref{sec:intro_aim}.

\subsection{Retrieval-Augmented Generation} \label{sec:back_rag}

\citet{lewis2020rag} introduced retrieval-augmented generation (RAG), which combines a pre-trained sequence-to-sequence model (its \emph{parametric} memory) with a dense vector index of Wikipedia that a neural retriever searches at query time (its \emph{non-parametric} memory). Generating from retrieved passages set new results on open-domain question answering and produced more specific and factual text than a model using only its parameters. Two properties of this design matter for a university assistant. First, the knowledge can be updated by replacing the document index rather than retraining the model, which suits policies and fee tables that change every year. Second, each answer can be traced back to the passages it was generated from, which makes citation possible.

Since then RAG has grown into a family of designs. \citet{gao2023rag} group them into \emph{naive} RAG (retrieve, then generate), \emph{advanced} RAG (adding steps such as query rewriting and re-ranking) and \emph{modular} RAG (flexible pipelines built from interchangeable parts). The survey also shows that most of the design decisions, such as how documents are split into chunks, how many passages are retrieved and which retriever is used, have no single best setting and must be chosen for the task. This motivates RQ2: rather than assuming a configuration, this project measures which choices matter on university documents.

The retriever itself can be lexical or dense. Lexical retrievers such as BM25 \cite{robertson2009bm25} score passages by the query words they contain, giving more weight to rare terms. They are fast, need no training and handle exact terms well, such as course codes, amounts and policy names, but they fail when the question and the document use different words. Dense retrievers encode questions and passages as vectors, either with a model trained for retrieval such as DPR \cite{karpukhin2020dpr} or with general sentence embeddings such as Sentence-BERT \cite{reimers2019sbert}, so that passages with the same meaning but different wording can still be matched. They can, however, miss rare exact terms. The two can be combined with reciprocal rank fusion \cite{cormack2009reciprocal}, which merges ranked lists using only the rank of each item and was shown to outperform more complex fusion methods. The pilot in Section~\ref{sec:feas_pilot} already shows the lexical weakness: a question about an ``extension'' failed to match a policy that only mentions ``extensions''. This is the reason for comparing lexical, dense and hybrid retrieval in RQ2.

\subsection{Hallucination and Grounding} \label{sec:back_halluc}

\citet{ji2023survey} define hallucination as generated content that is nonsensical or unfaithful to the provided source. They distinguish \emph{intrinsic} hallucination, which contradicts the source, from \emph{extrinsic} hallucination, which cannot be verified from it. Both matter here. An intrinsic error might state the wrong late-submission penalty; an extrinsic one might invent a refund rule that no document mentions. Retrieval reduces hallucination by giving the model the relevant text, but it does not remove it. If the retriever misses the right passage, the model may answer from its training data instead, which for a specific university is often outdated or simply wrong. For this reason the project measures faithfulness (whether each claim is supported by the retrieved passages) and citation accuracy separately from correctness, and treats an answer without support as a failure even if it happens to be true.

\subsection{Abstention and Confidence} \label{sec:back_abst}

A system that answers every question will eventually answer one it should not. \citet{wen2025know} survey \emph{abstention}, the ability of a language model to decline to answer, and analyse it from three perspectives: the query (for example, questions that are ambiguous or have no answer), the model (whether it has the knowledge and how confident it is) and human values (whether an answer would be harmful). They review methods across the model life cycle, from training to inference, and treat reliable abstention as an open problem at every stage. For a support assistant, abstaining is not enough on its own: the student still needs an answer, so the assistant should \emph{escalate} the question to staff.

Deciding when to escalate is a form of selective prediction, where a system may reject inputs on which it is likely to be wrong \cite{geifman2017selective}. The quality of such a system is described by a risk--coverage trade-off: answering fewer questions (lower coverage) should lower the error rate on the questions it does answer (lower risk). This framing gives the project a way to evaluate escalation as a measurable component, by plotting risk against coverage for different confidence thresholds rather than reporting a single accuracy figure. Possible confidence signals in RAG include the retrieval scores, the agreement between retrieved passages, and the model's own estimate of whether the passages support an answer; RQ3 tests whether such signals reduce incorrect answers without escalating too many answerable questions.

\subsection{University Support Chatbots} \label{sec:back_chat}

Chatbots for student support predate large language models. Early systems such as the NEU-chatbot for university admissions \cite{nguyen2021neu} used intent classification: each user question is mapped to one of a fixed set of intents, and a prepared answer is returned. This approach is predictable and cheap, but it only covers the questions its designers anticipated, and it can return a confident but wrong answer when a new question resembles an existing intent. The FAQ baseline (S1) in this project represents this approach. The pilot confirms its weakness: asked about a late-payment charge, it returned the answer for the fee due date.

More recent systems use LLMs. JayBot \cite{odede2024jaybot} is an LLM-based chatbot for students and admissions evaluated with users. Marcel \cite{trienes2025marcel} is the closest work to this project: a lightweight, open-source RAG assistant that answers prospective students' admission questions from university web pages, and adds a retriever that maps questions to curated FAQ entries so that staff can steer its answers. Marcel was evaluated on retrieval and generation using questions that the knowledge base can and cannot answer. However, its decision to abstain is left to the answer-generation prompt, and the authors note that other ways of judging answerability deserve investigation. Its scope is also narrow: admissions for prospective students rather than the wider range of questions current students ask, such as fees, regulations and mitigating circumstances.

\subsection{Evaluating RAG Systems} \label{sec:back_eval}

Evaluating a RAG system requires judging the retrieved passages and the generated answer, and doing so for many questions. RAGAs \cite{es2024ragas} proposes reference-free metrics computed with an LLM: faithfulness (whether the answer's claims are supported by the context), answer relevance and context relevance. LLM-based judging scales well, but \citet{zheng2023judging} show that LLM judges, while agreeing with human preferences at a level similar to agreement between humans, also have systematic biases, such as favouring the first answer shown, longer answers, or answers written by themselves. Automatic judgements must therefore be checked against human labels before they are trusted.

Unanswerable questions need their own evaluation. \citet{peng2025unanswerability} propose UAEval4RAG, a framework that generates unanswerable questions for a given knowledge base from a taxonomy of six categories and measures how often a RAG system correctly declines. They find that the choice of RAG components, such as the retriever and the prompt, strongly affects this behaviour. Most evaluations of university chatbots, by contrast, use small test sets and report averages without testing whether differences between systems are statistically meaningful.

\subsection{Summary and Research Gap} \label{sec:back_gap}

The literature shows that RAG can ground answers in documents, that grounding does not remove hallucination, that knowing when not to answer is a separate and hard problem, and that evaluation must cover unanswerable questions and be checked against humans. Three gaps follow for university support:
\begin{description}
  \item[G1] There is little controlled comparison of the main approaches (an intent or FAQ baseline, a plain LLM and RAG) on the same institutional knowledge base and test set. This leads to RQ1.
  \item[G2] Retrieval choices are usually fixed rather than measured on university documents, which mix web pages, policy PDFs and fee tables. This leads to RQ2.
  \item[G3] In existing systems such as Marcel, abstention is left to the prompt. It is not treated as a separate component, evaluated with a risk--coverage analysis, and linked to escalation to staff. This leads to RQ3.
\end{description}
This project addresses these gaps with a controlled comparison of four systems on Heriot-Watt University Malaysia documents, a test set that includes unanswerable questions, and an analysis of why each system succeeds or fails rather than a single accuracy score.

% Methodology
%%%%%%%%%%%%%%%%%%%%
\section{Methodology} \label{cha:method}
% TODO (Weeks 6-8): methods, requirements and research questions, linked to the aim and objectives.

\subsection{Research Approach} \label{sec:method_approach}

\subsection{The Four Systems} \label{sec:method_systems}

\subsection{Requirements} \label{sec:method_req}
% MoSCoW table of functional and non-functional requirements.

\subsection{Evaluation Plan} \label{sec:method_eval}
% Metrics, test set, development/test split, statistics.

\subsection{Success Criteria} \label{sec:method_success}

% Feasibility
%%%%%%%%%%%%%%%%%%%%
\section{Feasibility} \label{cha:feas}
% TODO (Weeks 5-8): design, evaluation strategy and preliminary work (pilot, document sources, tool tests).

\subsection{Document Sources} \label{sec:feas_docs}

\subsection{Pilot Study} \label{sec:feas_pilot}

\subsection{Tools and Technology Choices} \label{sec:feas_tools}

% Conclusions
%%%%%%%%%%%%%%%%%%%%
\section{Conclusions} \label{cha:concl}
% TODO (Week 10): short summary of the whole document.


%%%%%%%%%%%%%%%%%%%%%%%%%%%%%%%%%%%%%%%%%%%%%%%%%%
% BACK MATTER
%%%%%%%%%%%%%%%%%%%%%%%%%%%%%%%%%%%%%%%%%%%%%%%%%%

% References
% list of references

\newpage
\renewcommand{\bibname}{References} % change name to References
% \phantomsection
% \addcontentsline{toc}{section}{References}
\bibliography{references}
% \bibliographystyle{plainnat}

% using Harvard style
% \bibliographystyle{agsm}

% using ACM style (defined in acm_settings.tex)
\bibliographystyle{ACM-Reference-Format}

% using APA style (requires apacite package)
% \bibliographystyle{apacite}

% Appendices
\newpage
\appendix

\crefalias{chapter}{appendixchapter} % for appendix crossreferences

\section{Project Management}

\subsection{Gantt Chart}
% Semester 1 activities carried out and Semester 2 plan.

\subsection{Risk Analysis}
% Risk table with likelihood, impact, mitigation and contingency.

\subsection{Professional, Legal, Ethical and Social Issues}
% BCS code of conduct, document copyright and terms of use, data protection, ethics approval, social impact.

\section{Generative AI Summary}
% Required section: summary of GenAI use in the project and deliverables.

\section{Participant Information Sheet and Consent Form}
% Only if the usability study or peer labelling goes ahead (see the ethics application).

\section{Research Participant Questionnaire}
% Sample questionnaire, if the usability study goes ahead.

% Rest of appendices...

\end{document}
